%==============================================================================
\section{Thuật toán K-Nearest Neighbors (KNN)}
\label{sec:knn}
%==============================================================================

\begin{dinhnghia}[title={K-Nearest Neighbors (KNN)}]
K-nearest neighbors (KNN) là thuật toán học có giám sát dùng cho cả \textbf{phân loại} và \textbf{hồi quy}.
Trong công nghiệp, KNN chủ yếu dùng cho bài toán phân loại.
Ý tưởng: tìm $k$ điểm dữ liệu gần nhất với điểm kiểm tra và dùng các láng giềng đó để dự đoán.
$k$ là siêu tham số cần tinh chỉnh --- số láng giềng xét đến.
\end{dinhnghia}

\begin{phuongphap}[title={Phân loại}]
Với phân loại, KNN gán điểm kiểm tra vào \textbf{lớp xuất hiện nhiều nhất} trong $k$ láng giềng gần nhất (bỏ phiếu đa số).
\end{phuongphap}

\begin{phuongphap}[title={Hồi quy}]
Với hồi quy, KNN gán điểm kiểm tra bằng \textbf{trung bình} giá trị của $k$ láng giềng gần nhất.
\end{phuongphap}

\begin{dinhnghia}[title={Metric khoảng cách}]
Metric đo độ tương tự giữa hai điểm ảnh hưởng mạnh đến hiệu năng KNN.
Phổ biến: khoảng cách \textbf{Euclidean}, \textbf{Manhattan}, \textbf{Minkowski}.
\end{dinhnghia}

\begin{tinhchat}[title={Hai đặc tính định nghĩa KNN}]
\begin{itemize}
  \item \textbf{Lazy learning}: không có giai đoạn huấn luyện chuyên biệt; dùng toàn bộ dữ liệu khi phân loại.
  \item \textbf{Non-parametric}: không giả định gì về phân phối dữ liệu nền.
\end{itemize}
\end{tinhchat}

\subsection{KNN hoạt động như thế nào?}

\begin{ynghia}[title={Feature similarity}]
KNN dùng ``độ tương đồng feature'' để dự đoán giá trị điểm mới: điểm mới được gán nhãn dựa trên mức độ khớp với tập huấn luyện.
\end{ynghia}

\begin{phuongphap}[title={Bước 1 --- Dataset}]
Để triển khai bất kỳ thuật toán nào cần có dataset.
Ở bước đầu của KNN: nạp dữ liệu \textbf{huấn luyện} và \textbf{kiểm tra}.
\end{phuongphap}

\begin{phuongphap}[title={Bước 2 --- Chọn $K$}]
Chọn giá trị $K$ (số điểm láng giềng gần nhất). $K$ có thể là bất kỳ số nguyên nào.
\end{phuongphap}

\begin{phuongphap}[title={Bước 3 --- Với mỗi điểm trong tập test}]
\begin{enumerate}
  \item[\textbf{3.1}] Tính khoảng cách giữa điểm test và mỗi hàng dữ liệu huấn luyện (Euclidean, Manhattan hoặc Hamming; phổ biến nhất là Euclidean).
  \item[\textbf{3.2}] Sắp xếp theo khoảng cách tăng dần.
  \item[\textbf{3.3}] Chọn $K$ hàng đầu từ mảng đã sắp xếp.
  \item[\textbf{3.4}] Gán lớp cho điểm test theo \textbf{lớp xuất hiện nhiều nhất} trong $K$ hàng đó.
\end{enumerate}
\end{phuongphap}

\begin{phuongphap}[title={Bước 4 --- Kết thúc}]
Kết thúc quy trình cho mọi điểm test.
\end{phuongphap}

\subsection{Ví dụ minh họa}

\begin{vidu}[title={Phân loại điểm mới với $K = 3$}]
Giả sử có dataset có thể vẽ như sau (hai lớp màu):

\tsimg{03_knn/img01.jpg}

Cần phân loại điểm mới (chấm đen tại $(60, 60)$) vào lớp xanh hoặc đỏ.
Giả sử $K = 3$: tìm ba điểm gần nhất.

\tsimg{03_knn/img02.jpg}

Trong ba láng giềng, hai thuộc lớp Đỏ $\Rightarrow$ chấm đen được gán lớp \textbf{Đỏ}.
\end{vidu}

\subsection{Xây dựng mô hình KNN}

\begin{phuongphap}[title={Các bước xây mô hình}]
\begin{enumerate}
  \item \textbf{Nạp dữ liệu}: dùng pandas, numpy, hoặc \texttt{sklearn.datasets}.
  \item \textbf{Chia dữ liệu}: tập huấn luyện và kiểm tra.
  \item \textbf{Chuẩn hóa}: đảm bảo mỗi feature đóng góp như nhau vào khoảng cách.
  \item \textbf{Tính khoảng cách}: giữa điểm test và từng điểm huấn luyện.
  \item \textbf{Chọn $k$ láng giềng}: theo khoảng cách đã tính.
  \item \textbf{Dự đoán}: bỏ phiếu đa số (phân loại) hoặc trung bình (hồi quy).
  \item \textbf{Đánh giá}: accuracy, precision, recall, F1-score, v.v.
\end{enumerate}
\end{phuongphap}

\subsection{Triển khai KNN bằng Python}

\begin{ynghia}[title={Classifier và Regressor}]
KNN dùng được cho cả phân loại và hồi quy; dưới đây là hai ``recipe'' trong Python.
\end{ynghia}

\subsection{KNN làm bộ phân loại (Classifier)}

\begin{vidu}[title={Import và nạp Iris}]
\begin{lstlisting}
import numpy as np
import matplotlib.pyplot as plt
import pandas as pd
from sklearn import datasets

iris = datasets.load_iris()
headernames = ['sepal-length', 'sepal-width',
               'petal-length', 'petal-width', 'Class']
dataset = pd.DataFrame(
    np.column_stack([iris.data, iris.target_names[iris.target]]),
    columns=headernames)
dataset.head()
\end{lstlisting}

\textbf{Output mẫu \texttt{head()}:}
\begin{verbatim}
   sepal-length  sepal-width  petal-length  petal-width      Class
0           5.1          3.5           1.4          0.2  Iris-setosa
1           4.9          3.0           1.4          0.2  Iris-setosa
2           4.7          3.2           1.3          0.2  Iris-setosa
3           4.6          3.1           1.5          0.2  Iris-setosa
4           5.0          3.6           1.4          0.2  Iris-setosa
\end{verbatim}
\end{vidu}

\begin{vidu}[title={Tiền xử lý và chia tập}]
\begin{lstlisting}
X = dataset.iloc[:, :-1].values
y = dataset.iloc[:, 4].values

from sklearn.model_selection import train_test_split
X_train, X_test, y_train, y_test = train_test_split(
    X, y, test_size=0.40)

from sklearn.preprocessing import StandardScaler
scaler = StandardScaler()
scaler.fit(X_train)
X_train = scaler.transform(X_train)
X_test = scaler.transform(X_test)
\end{lstlisting}
\end{vidu}

\begin{vidu}[title={Huấn luyện và dự đoán}]
\begin{lstlisting}
from sklearn.neighbors import KNeighborsClassifier
classifier = KNeighborsClassifier(n_neighbors=8)
classifier.fit(X_train, y_train)
y_pred = classifier.predict(X_test)
\end{lstlisting}
\end{vidu}

\begin{vidu}[title={Báo cáo kết quả}]
\begin{lstlisting}
from sklearn.metrics import (classification_report,
    confusion_matrix, accuracy_score)

result = confusion_matrix(y_test, y_pred)
print("Confusion Matrix:")
print(result)
result1 = classification_report(y_test, y_pred)
print("Classification Report:")
print(result1)
result2 = accuracy_score(y_test, y_pred)
print("Accuracy:", result2)
\end{lstlisting}

\textbf{Output:}
\begin{verbatim}
Confusion Matrix:
[[21  0  0]
 [ 0 16  0]
 [ 0  7 16]]
Classification Report:
              precision    recall  f1-score   support
 Iris-setosa       1.00      1.00      1.00        21
Iris-versicolor       0.70      1.00      0.82        16
  Iris-virginica       1.00      0.70      0.82        23
     micro avg       0.88      0.88      0.88        60
     macro avg       0.90      0.90      0.88        60
  weighted avg       0.92      0.88      0.88        60

Accuracy: 0.8833333333333333
\end{verbatim}
\end{vidu}

\subsection{KNN làm bộ hồi quy (Regressor)}

\begin{vidu}[title={Dữ liệu hai feature, target petal-length}]
\begin{lstlisting}
import numpy as np
import pandas as pd
from sklearn import datasets

iris = datasets.load_iris()
headernames = ['sepal-length', 'sepal-width',
               'petal-length', 'petal-width', 'Class']
data = pd.DataFrame(
    np.column_stack([iris.data, iris.target_names[iris.target]]),
    columns=headernames)
array = data.values
X = array[:, :2]
Y = array[:, 2]
data.shape
\end{lstlisting}

\textbf{Output:} \texttt{(150, 5)}
\end{vidu}

\begin{vidu}[title={KNeighborsRegressor và MSE}]
\begin{lstlisting}
from sklearn.neighbors import KNeighborsRegressor

knnr = KNeighborsRegressor(n_neighbors=10)
knnr.fit(X, Y)
print("The MSE is:",
      format(np.power(Y - knnr.predict(X), 2).mean()))
\end{lstlisting}

\textbf{Output:}
\begin{verbatim}
The MSE is: 0.12226666666666669
\end{verbatim}
\end{vidu}

\subsection{Ưu và nhược điểm của KNN}

\begin{tinhchat}[title={Ưu điểm (Pros)}]
\begin{itemize}
  \item Thuật toán rất \textbf{đơn giản}, dễ hiểu và diễn giải.
  \item Hữu ích với dữ liệu \textbf{phi tuyến} vì không giả định dạng phân phối.
  \item \textbf{Đa năng}: dùng cho phân loại và hồi quy.
  \item Độ chính xác tương đối cao, dù còn nhiều mô hình có giám sát tốt hơn KNN.
\end{itemize}
\end{tinhchat}

\begin{luuy}[title={Nhược điểm (Cons)}]
\begin{itemize}
  \item \textbf{Tốn tính toán}: lưu toàn bộ dữ liệu huấn luyện.
  \item Cần \textbf{bộ nhớ lớn} so với thuật toán có giám sát khác.
  \item Dự đoán \textbf{chậm} khi $N$ lớn.
  \item \textbf{Nhạy thang đo} dữ liệu và feature không liên quan.
\end{itemize}
\end{luuy}

\subsection{Ứng dụng của KNN}

\begin{ynghia}[title={Ngân hàng}]
KNN dự đoán cá nhân có phù hợp để duyệt vay không: có đặc điểm giống nhóm vỡ nợ không?
\end{ynghia}

\begin{ynghia}[title={Xếp hạng tín dụng}]
So sánh với người có đặc điểm tương tự để ước lượng điểm tín dụng.
\end{ynghia}

\begin{ynghia}[title={Chính trị}]
Phân loại cử tri tiềm năng: ``Sẽ bầu'', ``Không bầu'', ``Bầu đảng X'', ``Bầu đảng Y'', v.v.
\end{ynghia}

\begin{ynghia}[title={Lĩnh vực khác}]
Nhận dạng giọng nói, chữ viết tay, ảnh và video.
\end{ynghia}
